\section{Schubert Calculus}

\subsection{The Chow Rings of Grassmanians}

We now return to our previous problem: given $V_1,\dots,V_4\subseteq\bP^{2n+1}$ four general $n$-planes,
how many projective lines intersect all of them.
We showed that this is just the size of the intersection
$$ Q = \bigcap\Sigma_n(V_i) \subseteq \bG(1,2n+1) . $$
To compute this we will study the Chow ring of $G=G(k,n)$, and we will show that the Schubert varieties
give a lot of information about the Chow ring (in fact, all of the information), and we can also compute
directly the products of Schubert varieties within the Chow group.
Let us keep this problem in mind while we discuss this.

We want to consider the classes $[\Sigma_a(\V)]\in\Ch(G(k,n))$.
Firstly we want to say two things: that the class does not depend on the complete flag $\V$, and that we
may multiply the classes by intersecting.
The following theorem gives us both.

\bthrm[name={kleiman}, title={Kleiman, special case}]

    Suppose $G=\gl_n(k)$ acts on $X$ algebraically.
    Then for any subvariety $Y\subseteq X$ and $A\in G$, $[AY]=[Y]$.

    Moreover if $G$ acts transitively on $X$ then for any two subvarieties $Y,Z\subseteq X$ then for
    general $A\in G$, $AX$ and $Y$ are generically transverse.

\ethrm

\bproof

    We only prove the first claim.
    Let $L$ be the line connecting $1$ to $A$ in $\Mat_n(k)$ and define
    $$ Z = \set{(B,x)\in(G\cap L)\times X}[B^{-1}x\in Y] , $$
    so that $Z$ is a subvariety of $\bP^1\times X$, and its fiber over $B=1$ is $Y$ and over $B=A$ is $AY$.
    Thus they are rationally equivalent as desired.
    \qed
    \TODO{Why is $Z$ is a variety}

\eproof


Consider a Schubert variety $\Sigma_a(\V)\subseteq G$.
Because $\gl_n$ acts algebraically and transitively on $G$, $[A\Sigma_a(\V)]=[\Sigma_a(\V)]$ by
\refmath*{kleiman}.
Because furthermore we have
$$ A\Sigma_a(\V) = \Sigma_a(A\V) , $$
and $\gl_n$ acts transitively on complete flags, we get that $[\Sigma_a(\V)]$ depends only on $a$ and not on
the flag $\V$.
Thus we define
$$ \sigma_a = [\Sigma_a(\V)] \in \Ch^{\abs a}(G) $$
for {\it any\/} flag $\V$.
These are called the {\emphcolor Schubert classes}.
And again by \refmath*{kleiman} we get that generally
$$ \sigma_{a_1}\cdots\sigma_{a_m} = \bracks{\bigcap_i\Sigma_{a_i}(\V_i)} . $$

Note that the $\Sigma_a(\V)$ form the closed strata of a stratification of $G$.
Recall that $\Sigma_b(\V)\subseteq\Sigma_a(\V)$ iff $a\leq b$ and thus the open strata are
$$ \Sigma^\circ_a = \Sigma_a - \bigcup_{a<b}\Sigma_b . $$
These are called {\emphcolor Schubert cells}.

\bprop

    The open strata $\Sigma^\circ_a$ are isomorphic to $\bA^{k(n-k)-\abs a}$.
    Thus the stratification is affine.

\eprop

\bnote

    The Schubert cells are important when the Schubert calculus is considered using the cohomology groups
    of $G(k,n)$ over $\bC$.

\enote

\bcoro

    The Chow group $\Ch(G)$ is generated by the Schubert classes.

\ecoro

\bprop

    In $\Ch(G)$ we have
    $$ \sigma_{n-k}^k = \sigma_{1^k}^{n-k} = \sigma_{(n-k)^k} \in \Ch_0(G) . $$

\eprop

\bproof

    The only Schubert class $\sigma_a$ where $\abs a=k(n-k)$ (i.e.\ $\sigma_a\in\Ch_0(G)$) is when
    $a=(n-k)^k$.
    Because $\sigma_a$ generate $\Ch(G)$ as an Abelian group, this means that $\sigma_{(n-k)^k}$ generates
    $\Ch_0(G)$.
    Because $G$ is a projective variety, we need only show that the degrees of $\sigma_{n-k}^k$ and
    $\sigma_{1^k}^{n-k}$ are one, and we will be done.

    By \refmath[example]{schubertinclusion} we have that for a codimension-$1$ subspace $H\subseteq V$,
    $$ \Sigma_{1^k}(H) = \set{\Lambda\leq V}[\Lambda\leq H] . $$
    Then for general codimension-$1$ subspaces $H_1,\dots,H_{n-k}$ we have that
    $$ \bigcap_i\Sigma_{1^k}(H_i) = \set{\Lambda\leq V}[\Lambda\leq\bigcap_iH_i] = \set{\bigcap_iH_i} . $$
    Furthermore it can be shown that this intersection is transverse, so
    $$ \sigma_{1^k}^{n-k} = \bracks{\bigcap_i\Sigma_{1^k}(H_i)} = \bracks{\set{\bigcap_iH_i}} $$
    has degree $1$.

    For $\sigma_{n-k}^k$ we can do the same argument or use the isomorphism $G(k,n)\cong G(n-k,n)$.
    \qed

\eproof

Because the Schubert classes generate $\Ch(G)$, we want to understand how they interact with one another.
We already know from Totaro [2014] that the $\sigma_a$ actually form a basis of $\Ch(G)$, but we will show
this directly.

\bdefn

    Two flags $\V,\W$ are {\emphcolor transverse} if any of the following equivalent conditions are true:
    \benum
        \item $V_i\cap W_{n-i}=0$ for all $i$,
        \item $\dim(V_i\cap W_j)=\min(0,i+j-n)$ for all $i,j$,
        \item There is a basis $e_1,\dots,e_n$ for $V$ such that
        $$ V_i = \gen{e_1,\dots,e_i},\qquad W_j = \gen{e_{n+1-j},\dots,e_n} . $$
    \eenum

\edefn

\bnote

    Clearly (3) implies (2) which implies (1).
    For the final implication let $e_i\in V_i\cap W_{n+1-i}$ be nonzero (which exists because
    $V_{i-1}\oplus W_{n+1-i}=V$ and thus the intersection has dimension $1$) and because
    $W_{n+1-i}\cap V_{n-1}$ is zero we have that $e_i\notin V_{i-1}$.
    This gives us $e_1,\dots,e_n\in V$ which are linearly independent and then $V_i=\gen{e_1,\dots,e_i}$.
    Furthermore $W_{n+1-i}=\gen{e_i,\dots,e_n}$.

\enote

We see that any two pairs of transverse flags can be carried from one to another by an element of $\gl_n$.
Furthermore we see that the set of pairs of transverse flags is open dense in the space of all pairs of flags.
From these two results we see that anything which is generally true for pairs of flags is true for
pairs of transverse flags and vice versa.
In particular if $\V,\W$ are transverse then $\Sigma_a(\V)\cap\Sigma_b(\W)$ are generically transverse.

\blemm

    Let $\V,\W$ be transverse flags, and let $\Lambda\in\Sigma_a(\V)\cap\Sigma_b(\W)$ be a general point
    of the intersection.
    \benum
        \item $\Lambda$ does not lie in any strictly smaller Schubert variety
        $\Sigma_{a'}(\V)\subsetneq\Sigma_a(\V)$.
        \item Define the induced flag $\Lambda^\V_i$ on $\Lambda$ by
        $$ \Lambda^\V_i = \Lambda \cap V_{n-k+i-a_i},\qquad i=1,\dots,k . $$
        Then $\Lambda^\V,\Lambda^\W$ are transverse.
    \eenum

\elemm

\bproof

    (1) holds by considering dimensions, as $\Sigma_{a'}(\V)\cap\Sigma_b(\W)$ has dimension strictly
    less than $\Sigma_a(\V)\cap\Sigma_b(\W)$.

    For (2) we want to show that
    $$ \Lambda^\V_i\cap\Lambda^\W_{k-i} = \Lambda\cap V_{n-k+i-a_i}\cap W_{n-i-b_{k-i}} = 0 $$
    for general $\Lambda$.
    \TODO{Complete?}

\eproof

\bprop

    Let $\V,\W$ be transverse flags, and $a,b$ be tuples with $\abs a+\abs b=k(n-k)$.
    Then $\Sigma_a(\V),\Sigma_b(\W)$ intersect transversely in a unique point if $b=a^*$ and are otherwise
    disjoint.
    Here we define $a^*$ to be
    $$ a^* = (n-k-a_k,\dots,n-k-a_1),\qquad a^*_i = n-k-a_{k+1-i} . $$
    Thus
    $$ \deg\sigma_a\sigma_b = \cases{1 & if $b=a^*$\cr0&else} . $$

\eprop

\bproof

    Because $\V,\W$ are transverse we have that the Schubert varieties intersect generically transversely.
    Thus
    $$ \deg\sigma_a\sigma_b = \abs{\Sigma_a(\V)\cap\Sigma_b(\W)} =
    \#\set{\Lambda}[\mltext{%
        $\dim(V_{n-k+i-a_i}\cap\Lambda)\geq i$,\cr
        $\dim(W_{n-k+i-b_i}\cap\Lambda)\geq i$%
    }\text{\quad for all $i$}] . $$
    If the intersection is nonzero, then let $\Lambda$ lie in the intersection.
    Then we have, for each $i$, the following
    $$ \dim(V_{n-k+i-a_i}\cap\Lambda) \geq i,\qquad \dim(W_{n-i+1-b_{k-i+1}}\cap\Lambda) \geq k+1-i . $$
    Thus we must have that $\Lambda\cap V_{n-k+i-a_i}\cap W_{n-i+1-b_{k-i+1}}$ is nonzero.
    Because the flags are transverse this means that
    $$ n - k + 1 - a_i - b_{k-i+1} \geq 1 \iff a_i + b_{k-i+1} \leq n - k . $$
    If this is an equality then $V_{n-k+i-a_i}\cap W_{n-i+1-b_{k-i+1}}$ meet in a one-dimensional subspace
    $\Gamma_i=\gen{e_{n-k+i-a_i}}$ necessarily contained in $\Lambda$.

    Now if $a_i+b_{k-i+1}\leq n-k$ then we have $\abs a+\abs b\leq k(n-k)$.
    Because this is an equality we must have $a_i+b_{k-i+1}=n-k$ for all $i$.
    Thus the disjoint $\Gamma_i$ are all contained in $\Lambda$, and thus $\Lambda=\bigoplus_i\Gamma_i$.
    So the intersection is a singleton if it is non-empty.
    \qed

\eproof

\bcoro

    The Schubert classes form a basis for $\Ch(G)$.

\ecoro

\bproof

    The Schubert classes generate $\Ch(G)$, and so if $\sigma_a=\sum n_c\sigma_c$ then we have
    $$ n_c = \deg(\sigma_a\sigma_{c^*}) , $$
    which is zero iff $a=c$.
    \qed

\eproof

Thus if we have a cycle of pure codimension $m$, we can compute its components directly.
That is if
$$ [\Gamma] = \sum_{\abs a=m}\gamma_a\sigma_a , $$
then
$$ \gamma_a = \deg([\Gamma]\sigma_{a^*}) = \#(\Gamma\cap\Sigma_{a^*}(\V)) $$
where $\V$ is a general flag (note that this makes sense only when $\Gamma\subseteq G$).
Thus
$$ [\Gamma] = \sum_{\abs a=m}\deg([\Gamma]\sigma_{a^*})\sigma_a . $$

In particular
$$ \sigma_a\sigma_b = \sum_{\abs c=\abs a+\abs b}\gamma_{a,b;c}\sigma_c,\qquad
\text{where } \gamma_{a,b;c} = \deg(\sigma_a\sigma_b\sigma_{c^*}) . $$
Here $\gamma_{a,b;c}$ are called the {\emphcolor Littlewood-Richardson coefficients}, and they depend only
on $a,b,c$ and not $n,k$.

Let $0\leq b\leq n-k$ be an integer, and define the {\emphcolor special Schubert class}
$\sigma_{(b)}=\sigma_b=\sigma_{b,0,\dots,0}$.
We drop the parentheses in $(b)$ when it is clear that $b$ is an integer and not a tuple.

\bthrm[title={Pieri's formula}, name={pieri}]

    For a Schubert class $\sigma_a\in\Ch(G)$ and a special Schubert class $\sigma_b$,
    $$ \sigma_b\sigma_a = \sum_{\mltext{$\abs c=\abs a+b$\cr$a_i\leq c_i\leq a_{i-1}$}}\sigma_c . $$

\ethrm

\bproof

    This is equivalent to showing that $\deg(\sigma_a\sigma_b\sigma_{c^*})$ is $1$ iff
    $a_i\leq c_i\leq a_{i-1}$ and zero otherwise.
    This can be shown by taking transverse flags $\V,\W$ for $a,c^*$ and a general flag $\U$ for $b$.
    \qed\TODO{This}

\eproof

\subsection{Returning to Our Problem}

Recall that we want to compute $\sigma_n^4$ in $\bG(1,2n+1)=G(2,2n+2)$.
So we focus here on $k=2$, and the dimension of $G$ is $4n$.
Using Pieri we have
$$ \sigma_n^2 = \sigma_{2n} + \sigma_{2n-1,1} + \cdots + \sigma_{n+1,n-1} + \sigma_{n,n} . $$
To square this again we want to compute products of Schubert classes of the form $\sigma_{a,b}$.

\blemm

    Assuming $a_1-a_2\geq b_1-b_2$,
    $$ \sigma_{a_1,a_2}\sigma_{b_1,b_2} = \sum_{k=0}^{b_1-b_2}\sigma_{a_1+b_1-k,a_2+b_2+k} . $$

\elemm

\bproof

    First consider the case where $b_1=b_2=b$, then by \refmath[example]{schubertinclusion}
    $$ \Sigma_{b,b}(\W) = \set{\Lambda}[\Lambda\leq W_{n-b}] . $$
    Thus
    $$ \eqalign{
        \Sigma_{a_1,a_2}(\V)\cap\Sigma_{b,b}(\W) &=
         \set{\Lambda}[\Lambda\cap V_{n-1-a_1}\neq0,\ \Lambda\subseteq V_{n-a_2},\ \Lambda\subseteq W_{n-b}]\cr
        &= \set{\Lambda}[\Lambda\cap(V_{n-1-a_1}\cap W_{n-b})\neq0,\ \Lambda\subseteq V_{n-a_2}\cap W_{n-b}]\cr
        &= \Sigma_{a_1+b,a_2+b}(\dots V_{n-1-a_1}\cap W_{n-b},\ V_{n-a_2}\cap W_{n-b}\dots) .
    } $$
    And so we have that
    $$ \sigma_{a_1,a_2}\sigma_{b,b} = \sigma_{a_1+b,a_2+b} . $$

    By Pieri we have that
    $$ \sigma_{a_1,a_2} = \sigma_{a_1-a_2}\sigma_{a_2,a_2} $$
    and when $b\leq a$
    $$ \sigma_a\sigma_b = \sum_{k=0}^b\sigma_{a+b-k,k} . $$
    And thus
    $$ \eqalign{
        \sigma_{a_1,a_2}\sigma_{b_1,b_2} &= \sigma_{a_1-a_2}\sigma_{b_1-b_2}\sigma_{a_2,a_2}\sigma_{b_2,b_2}\cr
            &= \sigma_{a_1-a_2}\sigma_{b_1-b_2}\sigma_{a_2+b_2,a_2+b_2}\cr
            &= \sum_{k=0}^{b_1-b_2}\sigma_{a_1-a_2+b_1-b_2-k,k}\sigma_{a_2+b_2,a_2+b_2}\cr
            &= \sum_{k=0}^{b_1-b_2}\sigma_{a_1+b_1-k,a_2+b_2+k} .
    } $$
    Which completes the proof.
    \qed

\eproof

This shows us that
$$ \sigma_{2n-k,k}^2 = \sum_{i=0}^{2n-2k}\sigma_{4n-2k-i,2k+i} . $$
Note that we must have $4n-2k-i\leq 2n$ so $i\geq 2n-2k$, thus
$$ \sigma_{2n-k,k}^2 = \sigma_{2n,2n} $$
which is the class of a point.
And when $j>k$,
$$ \sigma_{2n-k,k}\sigma_{2n-j,j} = \sum_{i=0}^{2n-2j}\sigma_{4n-k-j-i,k+j+i} , $$
so we have $4n-k-j-i\leq 2n$ so that $2n-k-j\leq i\leq 2n-2j$, but $j>k$ so this sum is empty.

Thus in
$$ \sigma_n^2 = \sum_k\sigma_{2n-k,k} , $$
all the terms square to the class of a point and distinct terms have product zero.
Thus
$$ \sigma_n^4 = (n+1)\sigma_{2n,2n} , $$
which means that $\deg\sigma_n^4=n+1$.
We have finally answered our original question.

\bansw

    If $V_1,\dots,V_4\subseteq\bP^{2n+1}$ are four general $n$-planes, then there are $n+1$ lines intersecting
    them all.

    In particular, for $n=1$, there are $2$ lines intersecting four general lines
    $L_1,\dots,L_4\subseteq\bP^3$.

\eansw

\subsection{A Generalization}

We now generalize the problem above, and show how we can use Schubert calculus to solve it once again.

\bprob

    Given four subvarieties $X_1,\dots,X_4\subseteq\bP^{2n+1}$ of dimension $n$ and degrees $d_1,\dots,d_4$,
    how many lines intersect all of them?

\eprob

First take a single subvariety $X$ of dimension $n$ and degree $d$.
Then consider the subvariety of lines intersecting $X$:
$$ \Sigma_X = \set{L\in\bG(1,2n+1)}[L\cap X\neq\emptyset] . $$
Let us further define the incidence subvariety
$$ I_X = \set{(x,L)\in X\times\bG(1,2n+1)}[x\in L] , $$
which comes equipped with two projections
{\def\diagrowheight{0pt}\diagram{
    &X\cr
    I_X\cr
    &\Sigma_X&\bG(1,2n+1)
}{
    \diagarrow{from={2,1}, to={1,2}, text={\pi_X}, y distance=.15cm}
    \diagarrow{from={2,1}, to={3,2}, text={\pi_\Sigma}, y distance=.15cm}
    \diagarrow{from={3,2}, to={3,3}, left cap={(}}
}}
Note that $\pi_X$'s fibers are isomorphic to $\bP^{2n}$ because the collection of lines in $\bP^{2n+1}$ which
pass through a specified point are isomorphic to $\bP^{2n}$.
Thus $\dim I_X=\dim X+2n=3n$.

Furthermore generally $\pi_\Sigma$ is one-to-one, and we can see that $\dim\Sigma_X=3n$ as well.
\TODO{Why?}

We know that within $\Ch(\bP^{2n+1})$, $[X]=d[\bP^n]$ by \refmath[corollary]{projvarinchow}.
Now consider $I=I_{\bP^{2n+1}}$, so that
$$ I_X = \pi_{\bP^{2n+1}}^{-1}X . $$
Pulling back along a morphism of varieties preserves rational equivalence (under certain conditions), and so
$$ [I_X] = \pi_\bP^*[X] = \pi_\bP^* d[\bP^n] = d[\pi_\bP^{-1}\bP^n] = d[I_{\bP^n}] \in \Ch(I) . $$
Now consider the projection $\pi_G\colon I\to\bG(1,2n+1)$, which pushes forward varieties:
$$ \pi_{G,*}[I_{\bP^n}] = [\pi_G(I_{\bP^n})] , $$
and
$$ \pi_G(I_{\bP^n}) = \set{L\in\bG(1,2n+1)}[L\cap\bP^n\neq\emptyset] = \Sigma_n(\bP^n) , $$
and thus
$$ \pi_{G,*}[I_{\bP^n}] = \sigma_n \in \Ch(G) . $$
Putting this together, we have that
$$ [\Sigma_X] = \pi_{G,*}[I_X] = d\pi_{G,*}[I_{\bP^n}] = d\sigma_n . $$

For general $X_i$ then we have, as computed before,
$$ \bracks{\bigcap_i\Sigma_{X_i}} = d_1\cdots d_4\sigma_n^4 = (n+1)d_1\cdots d_4\sigma_{2n,2n} . $$
The size of the intersection is the degree of this, which is $(n+1)d_1\cdots d_4$.

\bansw

    The number of lines intersecting four subvarieties $X_1,\dots,X_4\subseteq\bP^{2n+1}$ of dimensions $n$
    and degrees $d_1,\dots,d_4$ is $(n+1)d_1\cdots d_4$.

    In particular in $n=1$, the number of lines intersecting four curves $C_1,\dots,C_4$ of degrees
    $d_1,\dots,d_4$ is $2d_1\cdots d_4$.

\eansw

